% 第 0 章 Preliminaries (原书第 1-5 页)
\setcounter{chapter}{-1}
\chapter{预备知识}
\label{chap:0}

\section{引言}
\label{sec:0.1}
\index{Fourier 级数!收敛 (convergence of Fourier series)}

在本节中我们简要讨论本书的主题: Fourier 级数何时收敛? 我们令 $\mathbb{T}^n$ 表示 $n$ 维环面 $[0,1)^n$, 其元素为 $\theta=(\theta_1,\theta_2,\ldots,\theta_n)\in\mathbb{T}^n$, 而 $k=(k_1,k_2,\ldots,k_n)$ 为整数的 $n$ 元组. 若 $f\in L^p(\mathbb{T}^n)$, 则与 $f$ 相伴的 Fourier 变换和形式 Fourier 级数定义为
\begin{align*}
\hat{f}(k) &= \int f(\theta)e^{-2\pi ik\theta}\,d\theta, \\
f &= \sum_{k\in\mathbb{Z}^n}\hat{f}(k)e^{2\pi ik\theta}.
\end{align*}

本书的核心主题是分析 $f$ 的形式 Fourier 级数在何种意义下确实收敛于 $f$. 当然, 我们必须先取有限和, 再取极限. 由于求和取遍整数的所有 $n$ 元组, 即取遍所有 $k\in\mathbb{Z}^n$, 我们需要某种方法给 $\mathbb{Z}^n$ 的格点排序. 最简单的办法是把它们全部纳入不断扩张的球中: 定义 $|k|^2=k_1^2+k_2^2+\cdots+k_n^2$, 然后把 $f$ 的 Fourier 级数的第 $R$ 个球形部分和 (spherical partial sum) 定义为
\[
S_Rf(\theta)=\sum_{|k|<R}\hat{f}(k)e^{2\pi ik\theta}.
\]

在这个定义下, 我们要研究的基本问题是: $S_Rf$ 何时在 $L^p$ 中收敛到 $f$?

在一维情形, 答案是经典的, 由 M.~Riesz 于 $1910$ 年给出 (见 Zygmund~\cite{zygmund}, 第~7 章~2.4). 收敛对 $L^p$ 中所有的 $f$ 成立, 当且仅当 $1<p<\infty$.

在更高的维数, 这个问题直到最近才得到解决. Carl Herz 在 $1954$ 年证明, 收敛的一个必要条件是 $\frac{2n}{n+1}<p<\frac{2n}{n-1}$. Charles Fefferman 在 $1972$ 年证明, $L^p$ 中的收敛当且仅当 $p=2$ 时成立. Fefferman 的证明及其推论正是本书的主要内容.

Fefferman 的结果表明, 这不是给 Fourier 级数求和的正确方式. 那么还有别的方法吗? $\mathbb{Z}^n$ 的点本可以按球以外的某种次序分组. 同心多边形是一个显而易见的尝试, 但结果表明它并不比重复几个一维结果更有趣: 它不给出任何新数学, 还让人得以回避对 Fefferman 结果的深入思考. 回避思考一个课题几乎总是错误的; 往最好的方面说, 你日后也免不了要大吃几惊.

一位经典分析学家会立刻告诉你球形部分和的替代品: 使用求和法 (summability method). 我们将分析 S.~Bochner 引入的一种方法, 它本身是 Riesz 某个求和法的变形: 设 $(\xi)_+=\xi$ 当 $\xi>0$; 当 $\xi<0$ 时令它为 $0$. 于是我们定义 $\alpha$ 阶 \emph{Bochner--Riesz 平均} (Bochner--Riesz means) 为\index{Bochner--Riesz 平均 (Bochner--Riesz means)|emph}
\[
S_R^{\alpha}f(\theta)=\sum_{|k|<R}\left[\left(1-\frac{|k|^2}{R^2}\right)_{+}\right]^{\alpha}\hat{f}(k)e^{2\pi ik\theta}.
\]

关键在于: 若 $\alpha=0$, 则 $S_R^{\alpha}f=S_Rf$, 因此通过研究 $\alpha$ 趋于零时的极限行为, 我们有望理解 $S_R$ 的问题究竟出在哪里.

很好; 那么关于 $\alpha$ 接近零时的 $S_R^{\alpha}$, 已知的结果是什么? 高维时的行为极其复杂, 事实上完整的答案仅在二维时是已知的. 我们不打算写出答案的公式, 而是画一张图. 下面的图~\ref{fig:1} 与图~\ref{fig:2} 展示了 $S_R^{\alpha}$ 在 $1\leq p\leq\infty$ 上的 $L^p$ 有界性: 纵轴以 $\alpha$ 标度, 横轴以 $\frac{1}{p}$ 标度. 虚线与空心圆表示已知无界的点, 阴影区域表示已知有界的部分.

本书的目的, 就是给出图中各结果的详细证明.

\begin{figure}[htbp]
\centering
\begin{tikzpicture}[scale=0.9]
  % 阴影区域 (竖线填充)
  \fill[pattern=vertical lines, pattern color=black!70]
    (0,2) -- (1.25,0) -- (3.75,0) -- (5,2) -- (5,4) -- (0,4) -- cycle;
  % 横轴向外延伸的虚线
  \draw[dashed] (-0.5,0) -- (0,0);
  \draw[dashed] (5,0) -- (5.5,0);
  % 边框
  \draw (0,0) rectangle (5,4);
  % 虚线 V 形边界
  \draw[dashed] (0,2) -- (1.25,0);
  \draw[dashed] (3.75,0) -- (5,2);
  % 标记点
  \draw (0,2) circle (2.4pt);
  \draw (5,2) circle (2.4pt);
  \draw (1.25,0) circle (2.4pt);
  \draw (3.75,0) circle (2.4pt);
  \fill (2.5,0) circle (2.4pt);
  % 标注
  \node[left] at (-0.1,2) {$\alpha=\frac{1}{2}$};
  \node[below] at (1.25,-0.15) {$\frac{1}{4}$};
  \node[below] at (3.75,-0.15) {$\frac{3}{4}$};
\end{tikzpicture}
\caption{二维情形下 $S_R^\alpha$ 的 $L^p$ 有界性区域}
\label{fig:1}
\end{figure}

\begin{figure}[htbp]
\centering
\begin{tikzpicture}[scale=0.9]
  % 阴影区域 (竖线填充): 虚线 V 之上、实线 V 之上
  \fill[pattern=vertical lines, pattern color=black!70]
    (0,3.4) -- (1.5,0) -- (2.1,0.9) -- (2.6,0) -- (3.1,0.9) -- (3.7,0) -- (5,2.6) -- (5,4.6) -- (0,4.6) -- cycle;
  % 边框
  \draw (0,0) rectangle (5,4.6);
  % 虚线 V 形边界
  \draw[dashed] (0,3.4) -- (1.5,0);
  \draw[dashed] (3.7,0) -- (5,2.6);
  % 实线内边界
  \draw[thick] (2.1,0.9) -- (2.6,0) -- (3.1,0.9);
  % 标记点
  \draw (0,3.4) circle (2.4pt);
  \draw (5,2.6) circle (2.4pt);
  \draw (1.5,0) circle (2.4pt);
  \draw (3.7,0) circle (2.4pt);
  % 标注
  \node[left] at (-0.1,3.4) {$\alpha=\frac{n-1}{2}$};
  \node[left] at (-0.1,1.5) {$\alpha=\frac{n-1}{2(n+1)}$};
  \node[below] at (1.5,-0.15) {$\frac{n-1}{2n}$};
  \node[below] at (3.7,-0.15) {$\frac{n+1}{2n}$};
\end{tikzpicture}
\caption{高维情形下 $S_R^\alpha$ 的 $L^p$ 有界性区域}
\label{fig:2}
\end{figure}

\section{泛函分析的作用}
\label{sec:0.2}

本节给出攻克收敛问题的一个思路. 最好的办法是快速回顾一个标准结果: $\mathbb{T}$ 上存在连续函数, 其 Fourier 级数在 $\theta=0$ 处发散. 当然, 这个函数可以显式地构造出来, 但在一般维数下, 显式的计算是难以驾驭的. 我们转而求助于泛函分析的方法.

抽象的问题是这样的. 若我取一个有限 Fourier 级数, 其 Fourier 级数的部分和只是在进行到某处后停止, 我便收回原来的函数. 于是, Fourier 级数的部分和在所谓的三角多项式上收敛. 三角多项式在连续函数中稠密 (Stone--Weierstrass 定理), 而连续函数在 $L^p$ 中稠密. 我需要把``部分和在一个稠密子集上收敛''的信息过渡为``在整个函数空间上收敛'': 我需要交换极限. 标准的实分析告诉我, 某种一致性允许交换极限; 泛函分析则告诉我, 一致性是必要条件. 这正是共鸣定理 (uniform boundedness theorem) 的要义.

\begin{theorem}\label{thm:0.1}
\index{共鸣定理 (uniform boundedness theorem)}
存在 $f\in C(\mathbb{T})$, 使得 $\sup_N|S_Nf(0)|=\infty$.
\end{theorem}

\begin{proof}
我们先把这个问题转化为函数空间上的算子问题. 定义算子 $S_N:C(\mathbb{T})\to\mathbb{C}$ 为 $S_Nf=\sum_{|k|\leq N}\hat{f}(k)$. 由于和是有限的, 每个 $S_N$ 都是 $C(\mathbb{T})$ 上的连续线性泛函. 每个这样的线性泛函都由对一个有限 Borel 测度的积分给出; 在此情形,
\[
\sum_{|k|\leq N}\hat{f}(k)=\sum\int e^{-2\pi ik\theta}f(\theta)\,d\theta=\int K_N(\theta)f(\theta)\,d\theta,
\]
其中
\[
K_N(\theta)=\sum_{|k|\leq N}e^{-2\pi ik\theta}=\frac{\sin 2\pi\bigl(N+\frac{1}{2}\bigr)\theta}{\sin(\pi\theta)}.
\]
于是在此情形, 给出 $S_N$ 的测度是绝对连续的, 其全变差范数恰为 $\|K_N\|_1$. 这甚至不难计算; Zygmund~\cite{zygmund} 中给出了精确的计算:
\[
\|K_N\|_1=\frac{2}{\pi}\log N+O(1).
\]
真正的问题是把 $S_N$ 在个别函数上的作用, 与它作为空间 $C(\mathbb{T})$ 上的算子的行为联系起来. 这正是共鸣定理所告诉我们的: 要么 i)~$S_N$ 在算子范数下一致有界; 要么 ii)~对一个稠密集中的 $f$, $\sup_N|S_Nf|=\infty$. 由于第一个结论不成立, 第二个必然成立.
\end{proof}

这个结果绝佳地代表了本书余下部分要做的事情: 先用一些泛函分析的技巧, 把收敛问题化为函数空间上算子的有界性问题; 泛函分析之后, 是对具体算子的计算; 最终分析以一个详细的计算收尾.

\section{背景知识}
\label{sec:0.3}
\index{速降函数 (rapidly decreasing functions)}
\index{缓增分布 (tempered distributions)}

下一章将给出分析 Fourier 级数收敛所需的泛函分析. 最好的技巧是把问题从 Fourier 级数转移到 Fourier 积分. 这样做有好处: 显式计算一个积分, 比把一个级数求和成封闭形式要容易. 但也有坏处: 定义 Fourier 变换的那些积分并不绝对收敛. 克服这一困难的现代做法是, 先在一个万事皆收敛的稠密子集上工作, 然后再取极限. 这正是速降函数理论的要点, 本节对此作一概述. 更详细的论述见 Stein 与 Weiss~\cite{steinweiss}.

对 $f\in L^1(\mathbb{R}^n)$, Fourier 变换 $\hat{f}$ 定义为
\[
\hat{f}(\xi)=\int f(x)e^{-2\pi ix\xi}\,dx.
\]
该积分绝对收敛, 且 $\|\hat{f}\|_{\infty}\leq\|f\|_1$. 而要对 $f\in L^2$ 定义 $\hat{f}$, 就得耍些花招、用点骗术了.

速降函数空间或 Schwartz 函数空间记作 $\mathcal{S}$, 定义为 $\mathbb{R}^n$ 上所有使半范数 $\sup_x\allowbreak|x^{\alpha}D^{\beta}\phi(x)|$ 有限的光滑函数 $f$ 构成的类. 这里我们使用了多重指标记号: 若令 $\alpha$ 表示多重指标 $\alpha=(\alpha_1,\alpha_2,\ldots,\alpha_n)$, 则 $x^{\alpha}=x_1^{\alpha_1}\cdots x_n^{\alpha_n}$. 当这些半范数随 $\alpha$ 变动时, 它们在 $\mathcal{S}$ 上诱导出一个拓扑. $\mathcal{S}$ 的对偶空间称为缓增分布 (tempered distribution) 空间. 于是不难证明:

\begin{theorem}\label{thm:0.2}
\index{Fourier 变换!在 $\mathcal{S}$ 上 (Fourier transform on $\mathcal{S}$)}
\begin{enumerate}[label=\alph*)]
\item 若 $1\leq p<\infty$, 则 $\mathcal{S}$ 在 $L^p$ 中稠密.
\item Fourier 变换是 $\mathcal{S}$ 到 $\mathcal{S}$ 的连续单射满映射.
\item 若 $\mathcal{S}$ 中两个函数 $f$ 与 $g$ 的卷积定义为
\[
f\star g=\int f(x-y)g(y)\,dy,
\]
则该函数仍在 $\mathcal{S}$ 中, 并且 $\widehat{f\star g}=\hat{f}\,\hat{g}$.
\item 若 $f$ 与 $g$ 属于 $\mathcal{S}$, 则 $\int f\bar{g}=\int\hat{f}\,\overline{\hat{g}}$.
\item $\|f\|_2=\|\hat{f}\|_2$.
\item 分布 $u$ 的 Fourier 变换定义为 $\hat{u}(f)=u(\hat{f})$. 则 $\hat{u}$ 仍是一个分布.
\item 平移算子 $\tau_y$ 定义为 $(\tau_y f)(x)=f(x-y)$. 若把分布 $u$ 与 $\mathcal{S}$ 中函数的卷积定义为 $(u\star f)(x)=u(\tau_x\tilde{f})$, 其中 $\tilde{f}(y)=f(-y)$, 则 $u\star f$ 是光滑函数.
\end{enumerate}
\end{theorem}

空间 $\mathcal{S}$ 是一个技术性工具, 它使得我们能够方便地对那些原本可能发散的积分进行计算. 本书的其余部分还需要一些记号. 若 $T$ 是从 $L^p$ 到 $L^p$ 的有界线性算子, 我们把它的算子范数记作 ${}_p\|T\|$. 回忆
\[
{}_p\|T\|=\sup\{\|Tf\|_p \mid \|f\|_p=1\},
\]
以及
\[
{}_p\|T\|=\sup\left\{\left|\int Tfg\,\right| \mid \|f\|_p=\|g\|_{p'}=1\right\},
\]
其中 $\frac{1}{p}+\frac{1}{p'}=1$.

\section{第\texorpdfstring{$0$}{0}章评注}
\label{sec:0.4}
\index{Lebesgue 常数 (Lebesgue constants)}
\index{球形求和 (spherical summability)}

\noindent (0.1):\quad Bochner--Riesz 平均是 Riesz 平均 $(1-|n|^2)_+^{\alpha}$ 的变形; 我们将会看到, 这两种求和法在 $L^p$ 中给出相同的结果, 但 Bochner--Riesz 平均要容易计算得多. 对 $S_R^{\alpha}$ 最初的兴趣来自它与逐点求和及 Fourier 级数局部化问题的联系 (这个课题如今因其困难而名声不佳). Bochner 的论文~\cite{bochner} 表明, $n$ 维问题比一维问题更为深刻: 局部化并非对所有的 $\alpha$ 都成立, 这与一维的结果形成对照. 其中的基本思想也可见 Stein 与 Weiss~\cite{steinweiss} 第~7 章定理~4.2.

把 Fourier 系数按同心球而不是同心多边形分组, 是有充分理由的. 用 Bochner 自己的话说 (\cite{bochner}, 第~179--180 页):

\begin{quote}
``初等指数函数 $u(x)=e^{i(n_1x_1+\cdots+n_kx_k)}$ ($n_1,\ldots,n_k$ 均为整数) 是特征值问题
\[
\triangle u(x)=-\lambda u(x)
\]
的一组完备的正则解, 只要把这个系统视为定义在 (闭) 环面
\[
0\leq x_1\leq 2\pi,\ \ldots,\ 0\leq x_k\leq 2\pi
\]
上, 且 $\Delta$ 是关于环面上 Euclid 度量的 Laplace 算子, 即
\[
\frac{\partial^2 u}{\partial x_1^2}+\cdots+\frac{\partial^2 u}{\partial x_k^2}.
\]
由于 $\lambda=n_1^2+\cdots+n_k^2$, 我们书写级数的方式……满足一个十分自然的原则: 按特征值 $\lambda$ 的大小给级数的项排序.''
\end{quote}

用 Bochner 的话来复述: 在考察球形求和法时, 我们实际上是在考察紧 Riemann 流形上某个微分算子的特征函数展开的收敛性问题. C.~Fefferman 的结果已被用来证明: 当 $n>1$ 时, 这样的展开永不收敛. 这正是我们的主要论点: 球形平均的发散并非选择了古怪求和法所造成的反常, 而是高维几何的一个基本事实.

图~\ref{fig:1} 与图~\ref{fig:2} 所总结的结果归功于许多人. 概述可参见 Herz~\cite{herz}、C.~Fefferman~\cite{fefferman-spherical}、Carleson 与 Sjölin~\cite{christ}\footnote{译者注: 原书此处引作 Carleson and Sjölin [5], 但原书参考文献 [5] 实为 M.~Christ 的博士论文, 且原书参考文献缺失序号 6; 此处当指 L.~Carleson 与 P.~Sjölin 1972 年在 Studia Math. 上关于圆盘乘子问题的著名论文.} 以及 Tomas~\cite{tomas1}.

\noindent (0.2):\quad $\|K_R\|_1$ (称为 Lebesgue 常数) 的计算归功于 Fejér ($1910$ 年). 这类技巧在高维时一败涂地, 因为 $K_R(\theta)$ 的级数无法求和. 可以证明, 当 $n\geq 2$ 时 $\|K_R\|_1=O\bigl(R^{\frac{n-1}{2}}\bigr)$, 这已是最佳估计; 证明及更早结果的文献见 Shapiro~\cite{shapiro}. 在试图为部分和 $K_R$ 寻找显式公式时, 你很快会意识到: 一个真正好的公式将使你能够精确算出 $\mathbb{R}^n$ 中半径为 $R$ 的球内整点的个数. 不幸的是, 对这个量给出好的估计是数论中最困难的问题之一. 这是个严肃的问题, 因为对某些 $R$, 不存在整数对 $(m,n)$ 使 $m^2+n^2=R^2$; 而对稍微不同的 $R$, 这样的整数对又比比皆是. 你不能指望从 $K_R$ 的一个好公式得到任何真正规整的表达式.

关于一点处 Fourier 级数发散的经典详细证明, 最好查阅 Zygmund~\cite{zygmund} 第~7 章第一卷.
